\section{Prompts, Scores, and Notation}
\label{app:protocol}

\subsection{A common question and answer endpoint}

At the first answer position, let $\ell_{\mathrm{yes}}$ and
$\ell_{\mathrm{no}}$ denote the maximum log probabilities among the valid
single-token case and leading-space variants of each answer.  The registered
score is $P(\mathrm{yes})=\sigma(\ell_{\mathrm{yes}}-\ell_{\mathrm{no}})$,
where $\sigma$ is the logistic sigmoid. The original Edema cell uses ``pulmonary
edema,'' and the closure experiment uses ``consolidation'' in the same
template. Taking the best valid token variant for each answer makes the
endpoint insensitive to which listed capitalization or leading-space form
receives the highest score. This endpoint measures the relative support for
yes versus no at the first answer position; it does not measure the accuracy
of a generated clinical report.


The shared interface makes the question--direction grid interpretable:
changing a clinical direction does not change the question being scored.
Differences across columns still reflect different clinical questions, so
the ownership test compares directions within each column before drawing
a conclusion about the named concept.

\subsection{Reading the intervention notation}

Table~\ref{tab:notation} collects the notation introduced in the article.
The distinction between question $q$ and direction $d$ is central: a large
$W_{q,d}$ shows that a direction moves an answer, while $O_q$ asks whether
the direction bearing that answer's clinical name moves it more than its
competitors. A positive ownership margin alone is not the full confirmation
rule; the uncertainty and non-semantic control comparisons must also pass.
Keeping these roles separate allows the following analyses to distinguish
readability, intervention strength, and semantic attribution.

\begin{table}[H]
\centering
\caption{Notation used in the manuscript.}
\label{tab:notation}
\resizebox{\ifdim\width>\linewidth\linewidth\else\width\fi}{!}{%
\begin{tabular}{ll}
\toprule
Symbol & Definition \\
\midrule
$q,d,c$ & Clinical question, steering direction, and probe concept, respectively. \\
$\hvec_{it}$ & Consumed visual representation for image $i$ and visual token $t$. \\
$R,\boldsymbol{s}$ & Fixed Gaussian projection and train-fitted featurewise scale. \\
$\hat\wvec_c$ & Unit probe-normal direction for concept $c$ after lifting to the consumed representation. \\
$\alpha$ & Signed intervention dose; $|\alpha|$ is the perturbation norm relative to each token norm. \\
$W_{q,d}$ & Paired mean change in question $q$'s $P(\mathrm{yes})$ under direction $d$. \\
$O_q$ & Ownership margin $W_{q,q}-\max_{d\neq q}W_{q,d}$; positive values favor the named direction. \\
$S$ & Real-label AUROC minus the mean AUROC of the 20 nuisance-type control tasks. \\
$\gamma_i$ & Positive-minus-negative paired representation displacement along the Consolidation normal. \\
\bottomrule
\end{tabular}
}
\end{table}

The notation separates three levels of evidence: $S$ describes label
availability, $W_{q,d}$ describes an intervention effect, and $O_q$ compares
the named direction with clinical alternatives. The remaining appendix
follows this progression while keeping the evaluation cohorts explicit.

\FloatBarrier

